% STATUS: written 2026-11-06 (English gate clean; sources read in full as recorded below).
% Sources read in full for the cited passages: Johnson 1969 (pp. 92--93), Graf 2022 (p. 115).
\section{Kähler differentials and the universal derivation}\label{sec:commutative:kahler}

Let all rings here be commutative with identity. We look now for the single
derivation that carries all the others, as the introduction promised
\cite[p.~93]{Johnson1969Kahler}.

This section recalls that construction. It also fixes the notation used in
the rest of the chapter, so that one symbol serves for one object throughout.

Let $A$ be a ring and $B$ an $A$-algebra. Johnson constructs the module of
Kähler differentials from the multiplication map of $B$
\cite[p.~93]{Johnson1969Kahler}. Let $\mu\colon B\otimes_A B \to B$ be the
$A$-algebra map with $\mu(b\otimes b') = bb'$. Let $I$ be its kernel.

The quotient $I/I^2$ is a module over $B = (B\otimes_A B)/I$. There is an
$A$-derivation of $B$ into that module, given by
\[
  b \longmapsto (b\otimes 1 - 1\otimes b) + I^2
\]
\cite[p.~93]{Johnson1969Kahler}. So the kernel of the multiplication already
carries a derivation, and the derivation measures the failure of the two
factors to commute. The next definition names both
objects~\cite[p.~93]{Johnson1969Kahler}.

\begin{definition}[{\cite[p.~93]{Johnson1969Kahler}}]\label{def:comm:kahler}
Let $A$ be a ring and $B$ an $A$-algebra. With $I$ the kernel of the
multiplication $\mu\colon B\otimes_A B\to B$, the \emph{module of Kähler
differentials} of $B$ over $A$ is the $B$-module
$\Omega_{B/A} = I/I^2$. The \emph{universal derivation} is the $A$-linear map
$d_{B/A}\colon B \to \Omega_{B/A}$ given by
$d_{B/A}(b) = (b\otimes 1 - 1\otimes b) + I^2$.
\end{definition}

The name \emph{universal derivation} comes from a universal property that
Johnson states. Every derivation into a $B$-module factors through
$d_{B/A}$~\cite[p.~92]{Johnson1969Kahler}. The factor is unique
\cite[p.~92]{Johnson1969Kahler}. The next theorem records that property in
full~\cite[p.~92]{Johnson1969Kahler}.

It is the form the rest of this chapter uses, since every other derivation is
read through it.

\begin{theorem}[{\cite[p.~92]{Johnson1969Kahler}}]\label{thm:comm:universal}
Let $A$ be a ring and $B$ an $A$-algebra. Let $M$ be a $B$-module and let
$D\colon B\to M$ be an $A$-derivation. Then there is a unique $B$-linear map
$\varphi\colon \Omega_{B/A} \to M$ with $D = \varphi\circ d_{B/A}$.
\end{theorem}

Take $M = B$ in Theorem~\ref{thm:comm:universal}. This identifies the
derivations of $B$ with the dual of $\Omega_{B/A}$
\cite[p.~92]{Johnson1969Kahler}. Composition with $d_{B/A}$ maps
$\operatorname{Hom}_B(\Omega_{B/A}, B)$ bijectively to
$\operatorname{Der}_A(B)$~\cite[p.~92]{Johnson1969Kahler}.

So questions about the derivations of a commutative algebra are questions
about one module~\cite[p.~92]{Johnson1969Kahler}. When
$\operatorname{Der}_A(B)$ is free, one asks whether the geometry of $B$ is
smooth~\cite[p.~270]{BiswasGurjarKolte2014}. That is the Zariski--Lipman
problem of Section~\ref{sec:commutative:zariski-lipman}.

Graf takes the same duality as the definition of the tangent sheaf
\cite[p.~115]{Graf2022LZPosChar}. The tangent sheaf of a variety $X$ is the
dual of the sheaf of Kähler differentials
\cite[p.~115]{Graf2022LZPosChar}.

Section~\ref{sec:commutative:zariski-lipman} starts from this duality. It asks
what one learns about the singularities when the module of derivations is
free~\cite[p.~270]{BiswasGurjarKolte2014}.

A $K$-$E$-derivation of $K[x]$ is a map of the form $I-\varphi$ for an endomorphism $\varphi$, and the images of ideals under derivations and under these $E$-derivations are compared~\cite{WlImagesOfIdeals2008}.

For a finitely generated commutative algebra over a field of positive characteristic, the Lie subalgebras of its derivation algebra formed by locally finite derivations are studied~\cite{WlLieAlgebrasOfLocallyFiniteDeri2019}.

The images of the derivations of a polynomial algebra are described, which is the image question for the polynomial-algebra derivations of the chapter~\cite{WlImagesOfDerivationsOfPolynomia2017}.
